% !PRINTED 252
\weberpage{252}
Aus
\[
-A_3\gamma_3 = x x_0 x_1 x_2\left(\frac{1}{x}+\frac{1}{x_0}+\frac{1}{x_1}+\frac{1}{x_2}\right)
\]
ergibt sich nach (6) für $A_3$ der konstante Wert 1. Um aber $A_2$
zu bestimmen, müssen wir in der Entwickelung noch um ein Glied
weiter gehen. Wir setzen in
\[
x_0^4 + \gamma_3 x_0 - 27 = -A_2 x_0^2
\]
für $x_0$ die Entwickelung
\[
x_0 = -q^{-\frac{1}{3}} + 6\,q^{\frac{1}{3}} + \cdots,
\]
und finden $A_2 = 18$, so daß also die gesuchte Gleichung lautet:
\weq{7}{x^4 + 18\,x^2 + \gamma_3\,x - 27 = 0.}

Ist $n$ ungerade und nicht durch 3 teilbar, so nehme man $c$ \gesperrt{durch
12 teilbar an.}

Es ist alsdann
\[
\lambda \equiv n, \quad \alpha \equiv 0, \quad \delta \equiv 0, \quad
\beta \equiv a\,\eth_2, \quad \gamma \equiv -\eth\,a_2 \pmod{12}
\]
und es wird
\[
r^2 = (-1)^{\frac{\beta-1}{2}} = (-1)^{\frac{n-1}{2}}(-1)^{\frac{a-a'}{2}}, % [sic] printed a' here, not a_2 as in the parallel formulas above
\]
also der Satz

4. \gesperrt{Ist $n$ relativ prim zu 6, $c$ durch 12 teilbar, so
sind die Größen}
\[
P^{2}_{c,\eth,a}\,\gamma_2(\omega)^{n-1}\,\gamma_3(\omega)^{\frac{n-1}{2}}
\]
\gesperrt{Wurzeln einer Multiplikatorgleichung erster Stufe.}

Da die Funktionen $P^2$ für keinen endlichen Wert von $j(\omega)$
unendlich oder Null werden, so schließen wir, daß die Koeffizienten
der Gleichung, deren Wurzeln die $P^2$ sind, \gesperrt{ganze rationale
Funktionen von $\gamma_2, \gamma_3$ sind und daß der letzte
Koeffizient eine Konstante ist.}

Ist $n$ eine Primzahl $p$, so sind die Wurzeln dieser Gleichung
\[
x = p\left(\frac{\eta(p\omega)}{\eta(\omega)}\right)^2, \quad
x_h = (-1)^{\frac{p-1}{2}}\left(\frac{\eta\left(\dfrac{12h+\omega}{p}\right)}{\eta(\omega)}\right)^2
\quad h=0,1,\dots,p-1
\]
und die Anfangsglieder der Entwickelungen sind folgende:
\[
x = p\,q^{\frac{p-1}{6}}\dots, \quad
x_h = (-1)^{\frac{p-1}{2}}e^{\frac{2\pi i h}{p}}q^{-\frac{p-1}{6p}}\dots,
\]
wodurch sich für den letzten Koeffizienten der Wert $(-1)^{\frac{p-1}{2}}p$
ergibt. Daraus, daß das erste Glied in der Entwickelung von
%CONT>>
